\begin{figure*}[t]
\centering
\resizebox{\textwidth}{!}{%
\begin{tikzpicture}[node distance=6mm]
% center: the Laplacian
\node[box, minimum width=22mm] (lap) {Graph Laplacian\\$\Lap=\Deg-\Adj$};

% three operators
\node[hot, above right=6mm and 22mm of lap] (gcn) {GCN (local)\\$\hat{\Adj}$, $K$ hops};
\node[hot, right=22mm of lap] (heat) {Heat (global)\\$e^{-t\Lap}$};
\node[hot, below right=6mm and 22mm of lap] (qw) {Quantum walk (global)\\$e^{-it\Lap}$};

\draw[flow] (lap) -- (gcn);
\draw[flow] (lap) -- (heat);
\draw[flow] (lap) -- (qw);

% spreading law annotations
\node[plain, right=20mm of gcn] (gcns) {polynomial,\\$K$-hop horizon};
\node[plain, right=20mm of heat] (heats) {diffusive\\$\Var\sim t$};
\node[plain, right=20mm of qw] (qws) {\textbf{ballistic}\\$\Var\sim t^2$};

\draw[flow] (gcn) -- (gcns);
\draw[flow] (heat) -- (heats);
\draw[flow] (qw) -- (qws);

% mini spreading sketches
\begin{scope}[shift={($(heats.east)+(14mm,0)$)}]
  \draw[steel, thick] plot[domain=-1.2:1.2, samples=40]
       (\x, {0.6*exp(-2*\x*\x)});
  \node[font=\scriptsize, text=ink, above] at (0,0.65) {heat};
\end{scope}
\begin{scope}[shift={($(qws.east)+(14mm,0)$)}]
  \draw[amber!80!black, thick] plot[domain=-1.2:1.2, samples=60]
       (\x, {0.45*(cos(deg(6*\x))*exp(-1.5*\x*\x))+0.15});
  \node[font=\scriptsize, text=ink, above] at (0,0.65) {quantum walk};
\end{scope}
\end{tikzpicture}}
\caption{The unifying view. One Laplacian, three propagation operators, two spreading
laws. The only fundamental difference between classical diffusion and the quantum walk
is diffusive ($\Var\sim t$) versus ballistic ($\Var\sim t^2$) spreading; everything
else in the network is shared, which makes the operators param-matched.}
\label{fig:operators}
\end{figure*}
